\documentclass[11pt,a4paper]{amsart}
\usepackage[T1]{fontenc}
\usepackage{amsmath,amssymb,amsthm,mathtools,enumitem,microtype}
\usepackage[margin=2.7cm]{geometry}
\usepackage{xcolor}
\definecolor{linkcol}{rgb}{0.10,0.20,0.50}
\usepackage[colorlinks=true,linkcolor=linkcol,citecolor=linkcol,urlcolor=linkcol]{hyperref}
\theoremstyle{plain}
\newtheorem{theorem}{Theorem}[section]
\newtheorem{proposition}[theorem]{Proposition}
\newtheorem{lemma}[theorem]{Lemma}
\newtheorem{corollary}[theorem]{Corollary}
\theoremstyle{definition}
\newtheorem{definition}[theorem]{Definition}
\newtheorem{example}[theorem]{Example}
\newtheorem{remark}[theorem]{Remark}
\newcommand{\RR}{\mathbb R}
\newcommand{\QQ}{\mathbb Q}
\newcommand{\CC}{\mathbb C}
\newcommand{\NN}{\mathbb N}
\newcommand{\one}{\mathbf 1}
\newcommand{\Irr}{\operatorname{Irr}}
\newcommand{\Hom}{\operatorname{Hom}}
\newcommand{\im}{\operatorname{im}}
\newcommand{\indm}{\operatorname{ind}_{-}}
\newcommand{\Rad}{\operatorname{rad}}
\newcommand{\Sh}{\operatorname{Sh}}
\newcommand{\Vect}{\operatorname{Vect}}
\newcommand{\Hess}{\operatorname{Hess}}
\newcommand{\Def}{\operatorname{Def}}
\newcommand{\Pcal}{\mathcal P}
\newcommand{\Mcal}{\mathcal M}
\begin{document}
\title[The Stokes--microlocal Witt law]{The polarized Stokes--microlocal Witt law}
\author{M.~Haufschild}
\address{Independent researcher}
\email{haufschildmichael@gmail.com}
\date{4 August 2026}
\subjclass[2020]{Primary 34M40, 32S60; Secondary 14D24, 15A18, 18F20, 16G20}
\keywords{Stokes structure, irregular Riemann--Hilbert correspondence, microsupport,
vanishing cycle, characteristic cycle, Rees construction, wild character variety, Witt index}

\begin{abstract}
For a finite Galois-stable Puiseux spectrum, highest coefficient difference determines mixed Hom
irregularity.  We prove that its threshold partitions form a canonical dendrogram and give an
exact coarea decomposition of the full polarized Stokes index.  Every threshold term is a positive
square; physical rank is the sole potentially negative channel.

For every nonzero leading coefficient, we classify the zero directions of the actual exponential
phase $\operatorname{Re}(ce^{-ik\theta})$: they are $2k$ simple analytic walls, canonically split
into two transverse-sign fibers of size $k$.  On their induced finite stratification we construct
the cellular Stokes category.  Microlocalization is a quotient functor, nearby-to-vanishing is an
exact quotient, and the characteristic cycle of a degree-$k$ exponential Hom germ has $2k$
oriented rays of unit multiplicity, exactly $k$ in the ordered decay orientation.
Galois-normalized directed microstalk mass equals mixed Hom irregularity.  Thus the coarea
decomposition is also a categorical microlocal index formula.

After physical-rank fusion over an arbitrary marked curve, all local channels remain positive and
topology changes only one common-rank square.  The global form has negative index plus nullity at
most one for every genus, puncture count, ramification pattern, and nonnegative multi-Rees
specialization.  Capacity completion gives the sharp finite/affine/Lorentz trichotomy.  Corona
reduction is an exact positive Schur complement; its residual is the quotient in an exact
boundary--total--relative-vanishing sequence, and successive eliminations are order-independent.
An explicit three-term deformation complex derives the smooth-reduction dimension and Euler
formulas by rank--nullity.  These results give a mixed, polarized refinement of scalar irregular
rigidity formulas.
\end{abstract}
\maketitle

\section{Introduction}

Ramified Stokes geometry contains two gluing mechanisms.  Locally, exponential factors split into
branches on a finite cover and are regrouped by formal monodromy.  Globally, local Stokes spaces
are fused after physical ranks have been identified.  Main result: both mechanisms obey one
quadratic conservation law.  All local deformation channels form a positive coarea energy;
every possible negative contribution factors through total physical rank.

The scalar ancestry matters.  Hiroe proves, under multiplicity-free, regularity, and spectral
curve hypotheses, that an irregular rigidity index equals an Euler characteristic
\cite{Hiroe2019}; related local formulas involve Milnor numbers and intersection multiplicities
\cite{Hiroe2017}.  We do not present that scalar identity as new.  Our invariant is the entire
polarization on independent orbit charges.  Its coefficients are mixed $\Irr\Hom$ terms, and its
conclusion is an inertia law rather than one numerical evaluation.

Let $w^N=z$ be a common local cover and let
\[
                         q_a(w)=\sum_{k=1}^{K}c_{a,k}w^{-k}
\]
range through a finite $\mu_N$-stable set $A$ of distinct exponential factors.  For $a\ne b$,
let $k(a,b)$ be the largest exponent at which the coefficients differ and put
$h(a,b)=k(a,b)/N$; set $h(a,a)=0$.  Highest difference satisfies
\[
                         h(a,b)\leq\max\{h(a,c),h(c,b)\}.
\]
Hence threshold partitions form a canonical dendrogram.  Choose $H$ above all contacts.  If
$A_e$ are its blocks, $\ell_e>0$ their persistence lengths, $I=A/\mu_N$, and $y_a=x_{[a]}$, then
\begin{equation}\label{eq:intro-law}
 \boxed{Q_A(x)=\sum_{i\in I}x_i^2+
        \sum_e\ell_e\left(\sum_{a\in A_e}y_a\right)^2
        -(H-1)\left(\sum_{a\in A}y_a\right)^2.}
\end{equation}
No tree, cut family, or contact matrix appears among the hypotheses.  Branch blocks need not be
unions of Galois orbits.  Their exact descended vector is
$u_{e,i}=|A_e\cap i|$; assuming orbit saturation would lose ramified examples.

Formula \eqref{eq:intro-law} has four other readings.  Its positive linear forms are principal
conormals of cluster hyperplanes.  Their normalized coefficients are directed quotient
microstalks of exponential Hom germs.  They are associated-graded coordinates of a multi-Rees
family.  Their exact elimination is a corona Schur complement whose residual is a
relative-vanishing quotient.

For a genus-$g$ curve with $m\geq1$ marked points, identify local physical ranks and denote the
common value by $M$.  Define
\begin{equation}\label{eq:intro-fusion}
                      Q_{X,D}=\sum_{a=1}^{m}Q_{A_a}+2(1-g-m)M^2.
\end{equation}
Substituting \eqref{eq:intro-law} gives a positive-definite form plus one multiple of $M^2$.
Consequently
\begin{equation}\label{eq:intro-inertia}
 Q_{X,D}>0\ \text{on }\ker M\setminus\{0\},\qquad
 \indm(Q_{X,D})+\dim\Rad(Q_{X,D})\leq1.
\end{equation}
Here $\Rad(Q)$ is the radical of the polar bilinear form of $Q$; its dimension is the nullity of
$Q$.  Thus \eqref{eq:intro-inertia} is one joint bound, not two independent estimates.  This is
the global Stokes--microlocal Witt law.  Capacity completion replaces the auxiliary rank
coefficient by the intrinsic constrained minimum at $M=1$ and yields an exact
finite/affine/Lorentz trichotomy.  Every handle changes only this coefficient, by $-2$.

The argument uses finite coefficient, cellular Stokes, and tangent-reduction objects.  Analytic
irregular Riemann--Hilbert theory supplies surrounding realizations \cite{Sabbah,DKRH}; it is not
a proof interface.  We prove the exact transport statement: any exact realization equivalence
identifying the local incidence diagram preserves the quotient microstalk and all terms in (W).

Sections~\ref{sec:coefficients}--\ref{sec:global} prove coefficient and global laws.
Sections~\ref{sec:cellular}--\ref{sec:transport} construct categorical microlocal realization.
Sections~\ref{sec:rees}--\ref{sec:corona} treat Rees, capacity, and corona.
Section~\ref{sec:deformation} derives reduction dimension.  Final sections give marking legs,
spectral consequences, novelty boundaries, and falsifiers.

\section{Finite Puiseux spectra and canonical cuts}\label{sec:coefficients}

\begin{definition}[Equivariant finite spectrum]\label{def:spectrum}
Fix $N,K\geq1$ and a primitive $N$th root $\zeta$.  An equivariant finite Puiseux spectrum is a
finite set $A$, a permutation $\sigma:A\to A$ with $\sigma^N=1$, and coefficients
$c_{a,k}\in\CC$ such that
\[
                  c_{\sigma a,k}=\zeta^{-k}c_{a,k}
                  \qquad(a\in A,\ 1\leq k\leq K).
\]
Coefficient vectors are required to be pairwise distinct.
\end{definition}

For $a\ne b$, set
\[
 k(a,b)=\max\{k:c_{a,k}\ne c_{b,k}\},\qquad h(a,b)=k(a,b)/N,
\]
and put $k(a,a)=h(a,a)=0$.

\begin{lemma}[Highest-difference law]\label{lem:three-point}
For all $a,b,c\in A$,
\begin{equation}\label{eq:three-point}
                         k(a,b)\leq\max\{k(a,c),k(c,b)\}.
\end{equation}
Thus the largest of $k(a,b),k(a,c),k(b,c)$ occurs at least twice.
\end{lemma}

\begin{proof}
Above the maximum on the right, coefficients of $a,c$ agree and coefficients of $c,b$ agree.
Hence coefficients of $a,b$ agree.  Permuting the three variables excludes a unique largest value.
\end{proof}

For $t>0$, define
\begin{equation}\label{eq:threshold}
              a\sim_t b\quad\Longleftrightarrow\quad a=b\ \text{ or }\ h(a,b)<t.
\end{equation}

\begin{theorem}[Canonical coefficient dendrogram]\label{thm:dendrogram}
Every $\sim_t$ is an equivalence relation.  The resulting partitions $\Pcal_t$ coarsen with $t$.
Their finitely many blocks, ordered by inclusion, form a canonical rooted dendrogram.  For
distinct $a,b$, their smallest common block is born at $h(a,b)$.
\end{theorem}

\begin{proof}
If $a\sim_t b$ and $b\sim_t c$, Lemma~\ref{lem:three-point} gives $h(a,c)<t$; this proves
transitivity.  Monotonicity is immediate.  Only levels $k/N$ with $0\leq k\leq K$ occur.  Join
each block to the least strictly larger block and suppress repeated unary vertices.  Definition
\eqref{eq:threshold} identifies the birth level of the first common block.
\end{proof}

Choose $H>\max h(a,b)$.  Define the finite step integral
\begin{equation}\label{eq:partition-energy}
 E_H(y)=\int_0^H\sum_{C\in\Pcal_t}\left(\sum_{a\in C}y_a\right)^2dt,
\end{equation}
meaning interval length times the constant integrand between successive contact values.

\begin{theorem}[Coefficient coarea identity]\label{thm:coarea}
For every $y:A\to\RR$,
\begin{equation}\label{eq:coarea}
 E_H(y)=H\left(\sum_ay_a\right)^2-\sum_{a,b}h(a,b)y_ay_b.
\end{equation}
If $e$ runs through persistent non-root blocks $A_e$ and $\ell_e>0$ is persistence length, then
\begin{equation}\label{eq:cut-energy}
                         E_H(y)=\sum_e\ell_e\left(\sum_{a\in A_e}y_a\right)^2.
\end{equation}
Moreover $E_H(y)>0$ for $y\ne0$.
\end{theorem}

\begin{proof}
After expansion, the coefficient of $y_ay_b$ in \eqref{eq:partition-energy} is the length of
thresholds for which $a\sim_t b$, namely $H-h(a,b)$.  This proves \eqref{eq:coarea}.  Grouping the
step integral by persistent blocks gives \eqref{eq:cut-energy}.  At sufficiently small positive
$t$ all blocks are singletons, so a nonzero coordinate contributes on an interval of positive
length.
\end{proof}

\section{Orbit descent and the local Witt law}\label{sec:orbit}

Let $I=A/\langle\sigma\rangle$, write $[a]$ for the orbit of $a$, and put
\begin{equation}\label{eq:orbit-data}
 r_i=|i|,\qquad
 s_{ij}=\sum_{a\in i}\sum_{b\in j}h(a,b),\qquad
 M(x)=\sum_ir_ix_i.
\end{equation}

\begin{lemma}[Galois integrality]\label{lem:integrality}
Every $s_{ij}$ is a nonnegative integer.
\end{lemma}

\begin{proof}
Decompose $i\times j$ into cyclic pair orbits.  If one has length $L$ and leading difference
degree $k$, its stabilizer contains $\sigma^L$.  Equivariance gives $\zeta^{-Lk}=1$, so $N$ divides
$Lk$.  The pair orbit contributes $Lk/N\in\NN$.  Sum over pair orbits.
\end{proof}

\begin{definition}[Local polarized Stokes index]\label{def:local-index}
\begin{equation}\label{eq:local-index}
             Q_A(x)=\sum_i x_i^2+M(x)^2-\sum_{i,j}s_{ij}x_ix_j.
\end{equation}
\end{definition}

For a block $A_e$, define exact descended incidence and charge
\begin{equation}\label{eq:descended-cut}
        u_{e,i}=|A_e\cap i|,\qquad
        \nu_e(x)=\sum_i u_{e,i}x_i=\sum_{a\in A_e}x_{[a]}.
\end{equation}

\begin{theorem}[Local Stokes--microlocal Witt law]\label{thm:local-witt}
For every orbit charge $x$,
\begin{equation}\label{eq:local-witt}
 \boxed{Q_A(x)=\sum_i x_i^2+\sum_e\ell_e\nu_e(x)^2-(H-1)M(x)^2.}
\end{equation}
Hence $Q_A$ is positive definite on $\ker M$ and $\indm(Q_A)\leq1$.
\end{theorem}

\begin{proof}
Lift $x$ to $y_a=x_{[a]}$.  Then
\[
 \sum_{i,j}s_{ij}x_ix_j=\sum_{a,b}h(a,b)y_ay_b,
 \qquad \sum_ay_a=M(x).
\]
Insert Theorem~\ref{thm:coarea} into Definition~\ref{def:local-index}.  On $\ker M$, the final
square vanishes and $\sum_i x_i^2$ is strictly positive.  A two-dimensional negative subspace
would meet this codimension-one hyperplane nontrivially.
\end{proof}

\begin{remark}[Ramified descent]
Blocks $A_e$ may split a Galois orbit.  Formula \eqref{eq:descended-cut}, rather than an
orbit-saturation assumption, is required for ramified data.  When all blocks are orbit-saturated,
the usual weighted rooted-core cut vectors are recovered.
\end{remark}

\section{The global fused Witt law}\label{sec:global}

Let $D=\{p_1,\ldots,p_m\}$, $m\geq1$.  At $p_a$ choose a spectrum $A_a$, orbit space $I_a$,
local rank $M_a$, and local form $Q_a$.  Define
\begin{equation}\label{eq:fibre-product}
 V_{X,D}=\{(x_a)\in\bigoplus_a\RR^{I_a}:M_1(x_1)=\cdots=M_m(x_m)\}.
\end{equation}
Write $M$ for common rank and, for genus $g$, define
\begin{equation}\label{eq:fused-form}
             Q_{X,D}(x)=\sum_{a=1}^{m}Q_a(x_a)+2(1-g-m)M^2.
\end{equation}
Choose local root heights $H_a$ and set
\begin{align}\label{eq:global-positive}
 P_{X,D}(x)&=\sum_a\left(\sum_{i\in I_a}x_{a,i}^2+
                    \sum_e\ell_{a,e}\nu_{a,e}(x_a)^2\right),\\
 \Gamma_{X,D}&=2(1-g-m)-\sum_a(H_a-1).\notag
\end{align}

\begin{theorem}[Global coefficient-generated Witt law]\label{thm:global-witt}
On $V_{X,D}$,
\begin{equation}\label{eq:global-decomposition}
                  \boxed{Q_{X,D}=P_{X,D}+\Gamma_{X,D}M^2.}
\end{equation}
The form $P_{X,D}$ is positive definite, and
\begin{equation}\label{eq:global-inertia}
 Q_{X,D}>0\ \text{on }\ker M\setminus\{0\},\qquad
 \indm(Q_{X,D})+\dim\Rad(Q_{X,D})\leq1.
\end{equation}
\end{theorem}

\begin{proof}
Sum Theorem~\ref{thm:local-witt} and collect common-rank squares.  Terminal sums make
$P_{X,D}$ positive definite.  Any two-dimensional negative subspace meets $\ker M$ nontrivially.
The radical meets $\ker M$ trivially, so its dimension is at most one.  If a negative vector and a
nonzero radical vector coexisted, their span would contain a nonzero rank-zero vector; since the
radical is orthogonal to all vectors, that combination would have negative energy, contradiction.
\end{proof}

\begin{corollary}[Handle and sewing laws]\label{cor:handle}
Adding one handle subtracts $2M^2$ and changes no local channel.  Joining two marked curves away
from marked points adds their scalar rank coefficients and then subtracts $2$.
\end{corollary}

\begin{proof}
Use \eqref{eq:fused-form} and $\chi(X\#Y)=\chi(X)+\chi(Y)-2$.
\end{proof}

\section{Cellular Stokes categories and vanishing cycles}\label{sec:cellular}

Let $\mathsf C$ be the walking cospan $L\to W\leftarrow R$ and put
\[
                         \Sh_{\rm germ}=\operatorname{Fun}(\mathsf C,\Vect_{\QQ}).
\]
This is an abelian category.  For $S=L$ or $R$, write
$\rho_S(F):F(S)\to F(W)$ and define
\begin{equation}\label{eq:micro-functor}
                         \mu_S(F)=F(W)/\im\rho_S(F).
\end{equation}

\begin{proposition}[Microlocal quotient]\label{prop:microfunctor}
Equation \eqref{eq:micro-functor} defines a functor
$\mu_S:\Sh_{\rm germ}\to\Vect_{\QQ}$.  For every $F$,
\begin{equation}\label{eq:nearby-vanishing}
       F(S)\xrightarrow{\rho_S(F)}F(W)\longrightarrow\mu_S(F)\longrightarrow0
\end{equation}
is exact.
\end{proposition}

\begin{proof}
For a natural transformation $\alpha:F\to G$, naturality gives
$\alpha_W(\im\rho_S(F))\subseteq\im\rho_S(G)$.  Hence $\alpha_W$ induces a unique quotient map.
Identity and composition descend.  Exactness and surjectivity are the defining properties of a
quotient vector space.
\end{proof}

For a finite-dimensional vector space $V$ and $\varepsilon\in\{0,1\}$, define the exponential
germ $E(V,\varepsilon)$ by
\[
 E(W)=V,qquad (E(L),E(R))=
 \begin{cases}(V,0),&\varepsilon=1,\\(0,V),&\varepsilon=0,
 \end{cases}
\]
with identity and zero incidence maps.

\begin{corollary}[One oriented vanishing cycle]\label{cor:one-oriented}
If $V\ne0$, exactly one of $\mu_L(E(V,\varepsilon))$ and
$\mu_R(E(V,\varepsilon))$ is zero; the other is canonically $V$.
\end{corollary}

\begin{proof}
The cokernel of $V\xrightarrow{1}V$ is zero; the cokernel of $0\to V$ is $V$.
\end{proof}

Fix $c\in\CC^\times$ and $k>0$.  On the oriented boundary circle
$\RR/2\pi\mathbb Z$, use the dominance-wall convention
\begin{equation}\label{eq:analytic-phase}
 f_{c,k}(\theta)=\operatorname{Re}(ce^{-ik\theta})
  =|c|\cos(\arg c-k\theta).
\end{equation}
These are the analytic walls separating the two signs of exponential dominance.

\begin{proposition}[Complete analytic wall classification]\label{prop:analytic-walls}
Put $\tau_+=\pi/2$ and $\tau_-=-\pi/2$.  The zero set of \eqref{eq:analytic-phase} is
\begin{equation}\label{eq:analytic-wall-labels}
 \mathsf W_k(c)=\left\{
 \theta_{j,\varepsilon}=
 \frac{\arg c-\tau_\varepsilon}{k}+\frac{2\pi j}{k}\quad(\bmod\,2\pi):
 0\leq j<k,\ \varepsilon\in\{+,-\}\right\}.
\end{equation}
All $2k$ classes in \eqref{eq:analytic-wall-labels} are distinct and exhaust the zero set.  Every
wall is simple, with
\[
 f'_{c,k}(\theta_{j,+})=k|c|,
 \qquad f'_{c,k}(\theta_{j,-})=-k|c|.
\]
Consequently $\{0,\ldots,k-1\}\times\{+,-\}$ is canonically bijective with the analytic Stokes
walls, and its sign coordinate is their transverse orientation.
\end{proposition}

\begin{proof}
The equality $f_{c,k}(\theta)=0$ is equivalent on the angle circle to
$\arg c-k\theta=\pi/2$ or $-\pi/2$.  Multiplication by $k$ on
$\RR/2\pi\mathbb Z$ has, over each target, the $k$ distinct inverse images displayed in
\eqref{eq:analytic-wall-labels}.  The two fibers are disjoint because sine takes the values $1$
and $-1$ at their respective targets.  Finally
$f'_{c,k}(\theta)=k|c|\sin(\arg c-k\theta)$, which gives the asserted nonzero derivatives and
orientations.
\end{proof}

Use Proposition~\ref{prop:analytic-walls} to label the actual analytic walls by
$\mathsf W_k=\{0,\ldots,k-1\}\times\{+,-\}$.  Place one arc between consecutive walls.  Its
entrance quiver has two arrows from each arc to its endpoints.  No two nonidentity arrows compose,
so the path category contains only identities and incidence arrows.  Define
\begin{equation}\label{eq:boundary-category}
     \Sh_{\rm St}(S^1;k)=
     \operatorname{Fun}(\operatorname{Path}(\mathsf W_k),\Vect_{\QQ}).
\end{equation}
The alternating exponential object has wall stalk $V$, alternating arc stalks $V,0$, and identity
or zero incidence maps.

\begin{definition}[Cellular microsupport and cycle]\label{def:cellular-cycle}
An oriented conormal $(w,S)$ lies in cellular microsupport when
$\mu_S(F|_w)\ne0$.  Its cellular characteristic-cycle coefficient is
$\dim\mu_S(F|_w)$.  The directed cycle retains the $k$ conormals carrying ordered decay orientation.
\end{definition}

\begin{theorem}[Exponential microsupport]\label{thm:exponential-micro}
If $\dim V=d$, then the alternating degree-$k$ exponential object has exactly one microsupport
conormal over every wall, each nonzero cycle coefficient equals $d$, full cycle mass is $2kd$, and
directed decay mass is $kd$.
\end{theorem}

\begin{proof}
Each wall germ is $E(V,0)$ or $E(V,1)$, so the first two statements follow from
Corollary~\ref{cor:one-oriented}.  There are $2k$ walls and one directed wall in each of $k$
labelled antipodal pairs.
\end{proof}

Thus characteristic coefficients and vanishing ranks are dimensions of an actual quotient
functor, not names assigned to a numerical term.

\section{Microlocal Hom index and Galois descent}\label{sec:microlocal}

For branches $a,b$, let $E_{ab}$ be the exponential Hom object of degree $k(a,b)$.  If orbit $i$
carries a coefficient space $V_i$ of dimension $d_i$, tensor $E_{ab}$ by
$\Hom(V_{[a]},V_{[b]})$.

\begin{definition}[Normalized microlocal Hom mass]\label{def:micro-mass}
\begin{equation}\label{eq:micro-mass}
 \mathfrak m_{ij}(d)=\frac1N
 \sum_{a\in i}\sum_{b\in j}\sum_{\xi\in\operatorname{Decay}(a,b)}
       \dim\mu_\xi(E_{ab}\otimes\Hom(V_i,V_j)).
\end{equation}
\end{definition}

\begin{theorem}[Directed microlocal Hom index]\label{thm:micro-index}
For all orbit classes $i,j$,
\begin{equation}\label{eq:micro-index}
                         \mathfrak m_{ij}(d)=s_{ij}d_id_j.
\end{equation}
In particular, for one-dimensional coefficient spaces, normalized directed microstalk mass equals
mixed Hom irregularity.
\end{theorem}

\begin{proof}
Theorem~\ref{thm:exponential-micro} gives mass $k(a,b)d_id_j$ for one pair.  Sum and divide by
$N$ to obtain $s_{ij}d_id_j$.  Lemma~\ref{lem:integrality} proves integrality after normalization.
\end{proof}

Define $\langle x,x\rangle_\mu=\sum_{i,j}\mathfrak m_{ij}(1)x_ix_j$.

\begin{theorem}[Categorical microlocal coarea law]\label{thm:categorical-coarea}
\begin{align}\label{eq:categorical-coarea}
 Q_A(x)
 &=\sum_i x_i^2+M(x)^2-\langle x,x\rangle_\mu\notag\\
 &=\sum_i x_i^2+\sum_e\ell_e\nu_e(x)^2-(H-1)M(x)^2.
\end{align}
Hence cellular microsupport coefficients, directed vanishing-cycle ranks, mixed Hom irregularities,
and threshold coarea coefficients are four presentations of the same polarization.
\end{theorem}

\begin{proof}
The first equality is Definition~\ref{def:local-index} and Theorem~\ref{thm:micro-index}; the second
is Theorem~\ref{thm:local-witt}.
\end{proof}

\section{Exact realization transport}\label{sec:transport}

Let $\mathcal A$ be an abelian category and
$\Phi:\mathcal A\overset\sim\longrightarrow\Sh_{\rm germ}$ an exact equivalence.  Set
$\Psi_W=\operatorname{ev}_W\Phi$ and $\phi_S=\mu_S\Phi$.

\begin{theorem}[Irregular-realization transport]\label{thm:transport}
For every $A\in\mathcal A$,
\begin{equation}\label{eq:transport-exact}
 \operatorname{ev}_S\Phi(A)\longrightarrow\Psi_W(A)
 \longrightarrow\phi_S(A)\longrightarrow0
\end{equation}
is exact and functorial.  If $\Phi(A)$ is exponential, precisely one transported oriented
vanishing object has full coefficient rank.  Naturally isomorphic realization equivalences give
isomorphic quotient microstalks and preserve every term of \eqref{eq:categorical-coarea}.
\end{theorem}

\begin{proof}
Apply Proposition~\ref{prop:microfunctor} to $\Phi(A)$.  Exactness and functoriality survive
composition.  The exponential statement is Corollary~\ref{cor:one-oriented}.  A natural
isomorphism of realization functors identifies stalk maps and therefore their cokernels.
\end{proof}

\begin{remark}[Analytic scope]
Irregular Riemann--Hilbert theory realizes analytic Stokes objects in enhanced or Stokes-sheaf
categories \cite{Sabbah,DKRH}; enhanced nearby and vanishing cycles in dimension one are treated
in \cite{DKVanishing}.  Proposition~\ref{prop:analytic-walls} proves that the finite incidence
diagram is the stratification by the genuine leading exponential phase, not an abstract polygon.
Theorem~\ref{thm:transport} proves the implication used here once a realization identifies this
diagram.  We do not infer a meromorphic connection or a realization equivalence from the
quadratic form alone.
\end{remark}

\section{Principal conormals and the multi-Rees law}\label{sec:rees}

For each block $A_e$, let $\mathcal H_e=\{x:\nu_e(x)=0\}$.  Its constant principal conormal is
$d\nu_e=u_e=(u_{e,i})_i$.  Put
$P_A(x)=\sum_i x_i^2+\sum_e\ell_e\nu_e(x)^2$.

\begin{proposition}[Principal-conormal Hessian]\label{prop:hessian}
\begin{equation}\label{eq:hessian}
                   \frac12\Hess(P_A)=I+\sum_e\ell_e,d\nu_e\otimes d\nu_e.
\end{equation}
The sole subtractive Hessian term in $Q_A$ is $-(H-1)dM\otimes dM$.
\end{proposition}

\begin{proof}
Polarize the squares in Theorem~\ref{thm:local-witt}.
\end{proof}

For parameters $t_e\in\RR$, define
\begin{equation}\label{eq:local-rees}
 Q_{A,\boldsymbol t}(x)=\sum_i x_i^2+
       \sum_et_e\ell_e\nu_e(x)^2-(H-1)M(x)^2.
\end{equation}
The associated-graded coordinate of $e$ is $z_e=\sqrt{\ell_e}\nu_e$.  Unit specialization is
$Q_A$; setting $t_e=0$ deletes exactly channel $e$.

Choose independent $t_{a,e}$ at all marked points and fuse as in \eqref{eq:fused-form}.

\begin{theorem}[Global multi-Rees Witt law]\label{thm:global-rees}
If all $t_{a,e}\geq0$, the fused Rees form is a positive-definite terminal-plus-graded form plus
the same common-rank square as in Theorem~\ref{thm:global-witt}.  Thus
\[
 \indm(Q_{X,D,\boldsymbol t})+\dim\Rad(Q_{X,D,\boldsymbol t})\leq1,
 \qquad Q_{X,D,\boldsymbol t}>0\quad\text{on }\ker M\setminus\{0\}.
\]
The rank coefficient is independent of all Rees parameters.
\end{theorem}

\begin{proof}
Replace each local positive form by
$\sum_i x_i^2+\sum_et_{a,e}\ell_{a,e}\nu_{a,e}^2$.  Nonnegative parameters preserve all graded
squares and terminal sums remain positive definite.  Fusion changes only common rank.  Repeat the
hyperplane and radical argument in Theorem~\ref{thm:global-witt}.
\end{proof}

\section{Capacity completion and intrinsic phase}\label{sec:capacity}

Write $P_A(x)=x^TPx$ and $M(x)=r^Tx$, $r_i=|i|$.  Then
$Q_A(x)=x^TPx-\lambda(r^Tx)^2$, where $\lambda=H-1$ and $P>0$.  Define
\begin{equation}\label{eq:capacity-data}
                         c=r^TP^{-1}r>0,\qquad \Omega=c^{-1}-\lambda.
\end{equation}

\begin{theorem}[Capacity normal form]\label{thm:capacity}
For $m=r^Tx$, set $z=x-(m/c)P^{-1}r$.  Then $r^Tz=0$, this splitting is unique, and
\begin{equation}\label{eq:capacity-normal}
                         \boxed{Q_A(x)=z^TPz+\Omega m^2.}
\end{equation}
If $\Omega>0$, $Q_A$ is positive definite.  If $\Omega=0$, it is positive semidefinite with
kernel $\RR P^{-1}r$.  If $\Omega<0$, it is nonsingular with exactly one negative eigenvalue.
Moreover
\begin{equation}\label{eq:determinant}
                         \det Q_A=\det(P)c\Omega.
\end{equation}
\end{theorem}

\begin{proof}
The identity $r^TP^{-1}r=c$ gives $r^Tz=0$, uniqueness, and
$z^TP(P^{-1}r)=0$.  Expanding yields \eqref{eq:capacity-normal}; its orthogonal
positive-space-plus-one-line decomposition gives the trichotomy.  The determinant lemma gives
$\det(P-\lambda rr^T)=\det(P)(1-\lambda c)=\det(P)c\Omega$.
\end{proof}

\begin{proposition}[Root-height invariance]\label{prop:root-invariance}
Increasing $H$ by $t>0$ adds the full block $A$ with persistence $t$.  Hence
\[
 P_t=P+trr^T,\qquad \lambda_t=\lambda+t,
 \qquad c_t=\frac{c}{1+tc}.
\]
Thus $c_t^{-1}-\lambda_t=\Omega$ and $\det(P_t)c_t=\det(P)c$.
\end{proposition}

\begin{proof}
The full-block charge is $M=r^Tx$.  Sherman--Morrison gives the formula for $c_t$; the two
invariants follow, using the determinant lemma for the second.
\end{proof}

At puncture $a$, write $c_a,\Omega_a$ for local data and set $\delta_a=2-\Omega_a$.

\begin{theorem}[Euler--capacity refinement]\label{thm:euler-capacity}
Every $x\in V_{X,D}$ has unique centered components $z_a$, and
\begin{equation}\label{eq:euler-capacity}
 \boxed{Q_{X,D}(x)=\sum_a z_a^TP_az_a+\Omega_{X,D}M^2,
 \quad \Omega_{X,D}=\chi(X)+\sum_a(\Omega_a-2)
                    =\chi(X)-\sum_a\delta_a.}
\end{equation}
The coefficient $\Omega_{X,D}$ is the constrained minimum at $M=1$; its sign gives the global
positive, corank-one affine, and Lorentz phases.
\end{theorem}

\begin{proof}
Apply Theorem~\ref{thm:capacity} at every puncture and insert into \eqref{eq:fused-form}.  The rank
coefficient is
$\sum_a\Omega_a+2(1-g-m)=\chi(X)+\sum_a(\Omega_a-2)$.
The centered sum is positive definite and has unique minimum zero.
\end{proof}

\section{Corona elimination as relative vanishing}\label{sec:corona}

Let $X,Y$ be finite-dimensional real vector spaces and
\[
 \mathcal Q(x,y)=
 \begin{pmatrix}x\\y\end{pmatrix}^{T}
 \begin{pmatrix}A&B\\B^T&D\end{pmatrix}
 \begin{pmatrix}x\\y\end{pmatrix},\qquad D>0.
\]
Define
\begin{equation}\label{eq:corona-maps}
 h(x)=(x,-D^{-1}B^Tx),\qquad v(x,y)=y+D^{-1}B^Tx.
\end{equation}

\begin{theorem}[Exact corona--vanishing law]\label{thm:corona}
There is a short exact sequence
\begin{equation}\label{eq:corona-exact}
              0\longrightarrow X\xrightarrow{h}X\oplus Y
              \xrightarrow{v}Y\longrightarrow0.
\end{equation}
Moreover
\begin{equation}\label{eq:schur}
 \boxed{\mathcal Q(x,y)=v(x,y)^TDv(x,y)+x^T(A-BD^{-1}B^T)x.}
\end{equation}
Thus the Schur complement is the unique exact effective boundary energy.  Equality occurs exactly
on $\im h$, and the energy lost under elimination is the positive $D$-norm of the quotient
coordinate $(X\oplus Y)/\im h\cong Y$.
\end{theorem}

\begin{proof}
The map $h$ is injective, $v$ is onto since $v(0,y)=y$, and $v\circ h=0$.  Conversely,
$v(x,y)=0$ implies $(x,y)=h(x)$, proving exactness.  Expansion proves \eqref{eq:schur}; positivity
of $D$ proves the minimization and equality statements.
\end{proof}

\begin{corollary}[Order independence]\label{cor:order}
Successive eliminations of positive corona blocks are independent of elimination order.
\end{corollary}

\begin{proof}
Each step is exact minimization.  Iterated minima over a finite product commute.  Equivalently,
two attained lower bounds for the same full energy bound one another and are equal.
\end{proof}

The residual is not called a vanishing channel merely because its norm is positive: it is the
explicit quotient in \eqref{eq:corona-exact}.  Under a cellular realization,
Theorem~\ref{thm:transport} carries this quotient to the corresponding relative vanishing object.

\section{The Stokes deformation complex}\label{sec:deformation}

We derive the dimension identity from tangent data.  Fix $p$ punctures and integers
\[
 g,\quad G=\dim\mathfrak g,\quad z=\dim Z(\mathfrak g),\quad
 h_a,\ s_a,\ k_a\quad(1\leq a\leq p),
\]
where $h_a$ is formal-centralizer dimension, $s_a$ directed Stokes-root dimension, and $k_a$
prescribed conjugacy-class dimension.  Set
\begin{align}\label{eq:factor-counts}
 A_1&=2gG+\sum_a(G+h_a+s_a+k_a),\\
 R&=G-z+\sum_ah_a.\notag
\end{align}
At a smooth stable reduction point, let
\[
 d_{\rm act}:\CC^R\to\CC^{A_1},\qquad d_\mu:\CC^{A_1}\to\CC^R
\]
be effective gauge and moment derivatives.  Smooth stability here means $d_{\rm act}$ is
injective, $d_\mu$ is surjective, and $d_\mu d_{\rm act}=0$.

\begin{theorem}[Source-derived tangent dimension]\label{thm:tangent}
The reduced tangent $T=\ker d_\mu/\im d_{\rm act}$ satisfies
\begin{equation}\label{eq:tangent}
 \boxed{\dim T=2z+(2g-2+p)G+\sum_a(s_a+k_a-h_a).}
\end{equation}
\end{theorem}

\begin{proof}
Surjectivity gives $\dim\ker d_\mu=A_1-R$; injectivity gives
$\dim\im d_{\rm act}=R$.  Thus $\dim T=A_1-2R$.  Substitute \eqref{eq:factor-counts}.
\end{proof}

Form the three-term complex
\begin{equation}\label{eq:def-complex}
 \CC^z\oplus\CC^R\xrightarrow{d_0}\CC^{A_1}
 \xrightarrow{d_1}\CC^R\oplus\CC^z,
 \quad d_0(c,u)=d_{\rm act}u,\quad d_1(v)=(d_\mu v,0).
\end{equation}

\begin{theorem}[Self-dual Euler count]\label{thm:euler-count}
Sequence \eqref{eq:def-complex} is a complex and
$\dim H^0=z$, $H^1\cong T$, $\dim H^2=z$.  Hence
\begin{equation}\label{eq:euler-count}
 \boxed{\chi(\Def_{\rm St})=(2-2g-p)G+\sum_a(h_a-s_a-k_a).}
\end{equation}
\end{theorem}

\begin{proof}
The complex identity is $d_\mu d_{\rm act}=0$.  Injectivity leaves exactly the central summand in
$H^0$; degree-one cohomology is $T$; surjectivity leaves exactly the final central summand in
$H^2$.  Alternate dimensions and use Theorem~\ref{thm:tangent}.
\end{proof}

By Theorem~\ref{thm:micro-index}, $s_a$ is an integer directed microstalk sum.  Thus (M), (D), and
(E) give a closed chain from coefficient differences through vanishing quotients to tangent
dimension; a published dimension polynomial is not an input to this deduction.

\section{Type-A marking legs}\label{sec:legs}

Attach a path of length $l_i\geq0$ to orbit vertex $i$.  Put $y_{i,0}=x_i$,
$y_{i,l_i+1}=0$, and $d_{i,k}=y_{i,k}-y_{i,k+1}$.  Then $\sum_kd_{i,k}=x_i$.

\begin{lemma}[Path completion]\label{lem:path}
If $d_0+\cdots+d_l=x$, then
\begin{equation}\label{eq:path}
 \sum_{k=0}^{l}d_k^2=
 \sum_{k=0}^{l}\left(d_k-\frac{x}{l+1}\right)^2+\frac{x^2}{l+1}.
\end{equation}
\end{lemma}

\begin{proof}
Expand and use $\sum_kd_k=x$.
\end{proof}

\begin{theorem}[Leg-stable Witt law]\label{thm:legs}
Replacing each terminal square $x_i^2$ by its full type-$A$ path energy gives
\begin{align}\label{eq:leg-law}
 \widehat Q_A
 &=\sum_i\sum_{k=0}^{l_i}\left(d_{i,k}-\frac{x_i}{l_i+1}\right)^2
   +\sum_i\frac{x_i^2}{l_i+1}
   +\sum_e\ell_e\nu_e(x)^2-(H-1)M(x)^2.
\end{align}
Thus finite marking legs do not create a second negative or null direction.
\end{theorem}

\begin{proof}
Apply Lemma~\ref{lem:path} to every leg.  Fluctuation squares form a positive complement.  The
remaining positive matrix has diagonal $\operatorname{diag}((l_i+1)^{-1})>0$ plus the cut Gram
matrix and is perturbed by the same one rank square.  Repeat Theorem~\ref{thm:global-witt}.
\end{proof}

\section{Spectral consequence and sharpness}\label{sec:spectral}

Suppose an honest finite Stokes graph has symmetric adjacency matrix $B$ and its Cartan
polarization $C=2I-B$ is a form covered by Theorem~\ref{thm:global-witt}, including its marking
legs.  Let $\lambda_1(B)\geq\cdots\geq\lambda_n(B)$.

\begin{theorem}[Second-eigenvalue obstruction]\label{thm:lambda-two}
If $n\geq2$, then $\lambda_2(B)\leq2$.  If $\lambda_1(B)=2$, eigenvalue $2$ is simple.  If
$\lambda_1(B)>2$, exactly one eigenvalue exceeds $2$ and none equals $2$.
\end{theorem}

\begin{proof}
Eigenvalues of $C$ are $2-\lambda_i(B)$.  Apply \eqref{eq:global-inertia}.
\end{proof}

\begin{example}[Sharp Lorentz phase]\label{ex:triangle}
For
\[
 B=\begin{pmatrix}0&6&4\\6&0&3\\4&3&0\end{pmatrix},\qquad C=2I-B,
\]
one has $\one^TC\one=-20$, $\det C=-258$, and $\operatorname{tr}C=6$.  Hence inertia is
$(2,1,0)$.  One negative direction is sharp.
\end{example}

\section{Novelty boundary and falsifiers}\label{sec:novelty}

Finite Stokes filtrations and categorical forms are classical \cite{Sabbah}; irregular
Riemann--Hilbert theory supplies analytic realizations \cite{DKRH}; ultrametric matrices have tree
and inverse $M$-matrix theories \cite{Martinez1994,McDonald1995,Nabben1995}; fission and wild
character spaces are developed in \cite{Boalch2014}; Hiroe gives a scalar spectral-curve Euler
identity \cite{Hiroe2019}.  These are genuine ancestors.

New claim is their mixed dependency-closed polarization:
\begin{align*}
 \text{Puiseux coefficients}&\Longrightarrow\text{canonical cuts}
 \Longrightarrow\text{quotient microstalk masses}\\
 &\Longrightarrow\text{mixed Hom index}
 \Longrightarrow\text{positive coarea plus one rank line}.
\end{align*}
Global content is not a scalar rigidity identity: all punctures and handles share the same one
rank port, forcing the Witt bound in arbitrary genus.  Global multi-Rees stability and exact
corona relative-vanishing describe how collateral constructions preserve that port.

Any of the following would falsify the theorem as stated.
\begin{enumerate}[label=(\arabic*),leftmargin=2.2em]
\item an equivariant spectrum violating Lemma~\ref{lem:three-point};
\item failure of finite coarea identity \eqref{eq:coarea};
\item a zero of $\operatorname{Re}(ce^{-ik\theta})$ outside
      \eqref{eq:analytic-wall-labels}, a repeated wall, or failure of the transverse-sign law;
\item a directed exponential germ whose quotient microstalk mass differs from its degree;
\item a nonzero compatible rank-zero charge of nonpositive fused energy;
\item a nonnegative Rees specialization whose negative index plus nullity exceeds one;
\item a positive corona block violating \eqref{eq:schur} or \eqref{eq:corona-exact}.
\end{enumerate}
Conversely, this is not an analytic existence criterion: a quadratic form satisfying the Witt law
need not arise from a meromorphic connection.

\section{Conclusion}

The universal law factors the full mixed Stokes index through a positive microlocal coarea space
and one physical-rank line.  Principal conormals, cellular characteristic coefficients, vanishing
quotients, Rees graded channels, corona residuals, marking legs, and local terms of the deformation
complex are positive coordinates of that same space.  Fusion and topology cannot create another
defect because they see local data only through common physical rank.

This single structure explains analytic wall count and orientation, local Lorentz bounds, affine
simplicity, capacity walls, handle
monotonicity, Rees stability, Schur order independence, and the adjacency bound
$\lambda_2\leq2$.  It also replaces the former interface between coefficient contacts and
microlocal multiplicities by an exact categorical quotient calculation.

\begin{thebibliography}{99}
\bibitem{Boalch2014}
P.~Boalch, \emph{Geometry and braiding of Stokes data; fission and wild character varieties},
Ann. of Math. (2) \textbf{179} (2014), 301--365.
\bibitem{DKRH}
A.~D'Agnolo and M.~Kashiwara, \emph{Riemann--Hilbert correspondence for holonomic
$\mathcal D$-modules}, Publ. Math. Inst. Hautes Etudes Sci. \textbf{123} (2016), 69--197.
\bibitem{DKVanishing}
A.~D'Agnolo and M.~Kashiwara, \emph{Enhanced nearby and vanishing cycles in dimension one and
Fourier transform}, Publ. Res. Inst. Math. Sci. \textbf{59} (2023), 543--570.
\bibitem{Hiroe2017}
K.~Hiroe, \emph{Ramified irregular singularities of meromorphic connections and plane curve
singularities}, Josai Math. Monogr. \textbf{10} (2017), 161--192.
\bibitem{Hiroe2019}
K.~Hiroe, \emph{Index of rigidity of differential equations and Euler characteristic of their
spectral curves}, arXiv:1907.12846.
\bibitem{Martinez1994}
S.~Mart\'{i}nez, G.~Michon and J.~San Mart\'{i}n, \emph{Inverse of strictly ultrametric matrices
are of Stieltjes type}, SIAM J. Matrix Anal. Appl. \textbf{15} (1994), 98--106.
\bibitem{McDonald1995}
J.~J.~McDonald, M.~Neumann, H.~Schneider and M.~J.~Tsatsomeros, \emph{Inverse $M$-matrix
inequalities and generalized ultrametric matrices}, Linear Algebra Appl. \textbf{220} (1995),
321--341.
\bibitem{Nabben1995}
R.~Nabben and R.~S.~Varga, \emph{Generalized ultrametric matrices---a class of inverse
$M$-matrices}, Linear Algebra Appl. \textbf{220} (1995), 365--390.
\bibitem{Sabbah}
C.~Sabbah, \emph{Introduction to Stokes Structures}, Lecture Notes in Mathematics 2060,
Springer, 2013.
\end{thebibliography}
\end{document}
